\documentclass[9pt,twocolumn,twoside]{pnas-new}
% Use the lineno option to display guide line numbers if required.

\articletype{Social Sciences}%%%% Article topic/classification

\templatetype{pnasresearcharticle} % Choose template
% {pnasresearcharticle} = Template for a two-column research article
% {pnasmathematics} %= Template for a one-column mathematics article
% {pnasinvited} %= Template for a PNAS invited submission


\begin{document}


%%% solo jugando con los titulos, tiremos ideas %%%%
\title{The Rise and Fall of Automatic Behavior: Evidence from Discrete Choice Experiments with Eye Tracking}

% Use letters for affiliations, numbers to show equal authorship (if applicable) and to indicate the corresponding author
\author[a,c]{Omar David Perez\thanks{Corresponding author: Omar D. Perez, omar.perez.r@uchile.cl, 8370439 Santiago, Chile}}
\author[b,c]{Bastian Henriquez-Jara\thanks{Corresponding author: Bastián Henríquez-Jara, bastian.henriquez@uchile.cl, 8370439 Santiago, Chile}}
\author[d]{Sebastian Muñoz}
\affil[a]{Department of Industrial Engineering, Universidad de Chile, Santiago, Chile}
\affil[b]{Department of Civil Engineering, Universidad de Chile, Santiago, Chile}
\affil[c]{Complex Engineering Systems Institute (ISCI), Santiago, Chile}
\affil[d]{NeuroPolicy Lab, Universidad de Chile, Santiago, Chile}

% Please give the surname of the lead author for the running footer
\leadauthor{Perez}

% Please add a significance statement to explain the relevance of your work
\significancestatement{Public investments, mitigation plans and demand forecast analysis rely on Discrete Choice Experiments (DCEs) to validate people's preferences and potential social and private welfare implications. A well-established assumption behind this economic analysis is that people's choices in DCEs will reflect their underlying stable preferences, which are reflected as they are presented with a limited number of choice experiences. We show in this paper that this is not guaranteed. When people make repeated decisions for long periods in the same environment (such as when they go to the supermarket every week or take the bus everyday) they tend to develop different automatic habits which repeat previous choices, without directing their attentional resources to the relevant options' attributes. However, if these attentional resources are captured by changes in the environment stability, they will increase the attentional breadth and behave in a more rational, utility-maximizing manner, directing attention to relevant attributes. Behavior automaticity is a consequence of people's tendency to deploy their visual attention.} 

% Please include corresponding author, author contribution and author declaration information
\authorcontributions{O.D.P. and B.H.J designed the study. S.M. collected the data. S.M. and B.H.J. analyzed the data. O.D.P. and B.H.J. wrote the paper.}
\authordeclaration{Please declare any competing interests here.}
\equalauthors{\textsuperscript{1}O.D.P. contributed equally to this work with B.H.J.}
\correspondingauthor{\textsuperscript{2}To whom correspondence should be addressed. E-mail: omar.perez.r@uchile.cl}

% At least three keywords are required at submission. Please provide three to five keywords, separated by the pipe symbol.
\keywords{Discrete choice $|$ Attention $|$ Habits $|$ Inertia $|$ Eye-tracking}

\begin{abstract}
Discrete-choice experiments are widely employed for policy and welfare analysis as they are assumed to reflect people's stable, underlying preferences for specific options' attributes. However, constantly calculating optimal tradeoffs is cognitively demanding and may not persist when repeated choices are made in stable environments. Using eye-tracking and computational modeling in a DCE, we show that extended repetition in stable environments fundamentally changes the cognitive architecture of decision-making, where different behavioral strategies are deployed conditionally on how attentional resources are deployed. While some subjects maintain exploratory attention and behave as utility-maximizers, another group offloads attention from attributes, relying on spatial, habit-like heuristics. Critically, we show that choice inertia---the marker of automaticity---is endogenous to visual gaze entropy. By including changes in environmental stability, attention is captured and habitual subjects instantly break their heuristics. Contrary to widely held assumptions, joint computational modeling demonstrates that preferences are dynamically constructed and affected by the level of attentional load. Our results thus bridge cognitive mechanisms with structural econometrics to resolve a foundational problem in discrete-choice modeling. 
\end{abstract}

\dates{This manuscript was compiled on \today}
\doi{\url{www.pnas.org/cgi/doi/10.1073/pnas.XXXXXXXXXX}}


\maketitle
\thispagestyle{firststyle}
\ifthenelse{\boolean{shortarticle}}{\ifthenelse{\boolean{singlecolumn}}{\abscontentformatted}{\abscontent}}{}

\Firstpage

\dropcap{I}n economics, marketing, and transportation research, discrete choice experiments (DCEs) are a widely used method for eliciting preferences and estimating people's willingness-to-pay (WTP) for different attributes of goods and services---travel times in transportation or air quality in environmental studies, for example. DCEs are then used to predict how people will respond to potential changes in those attributes, making them a fundamental tool for policy and business decisions \cite{mcfadden1974, train2009}.

In a typical DCE, participants choose between two or more alternatives that differ in their attributes (\autoref{fig:framework}). These choices are used to estimate the parameters of a utility function representing their underlying, stable preferences. This assumption is critical because if preferences change across the task, the estimated parameters will not reflect the true preferences of the participants, and predictions based on those estimates will be inaccurate, making DCEs unreliable for policy and business decisions. A second, fundamental assumption of DCEs is that preferences are formed by considering and comparing the whole set of attributes of each alternative. If participants do not process all of the attributes, the estimated parameters would reflect only the subset of attributes that were attended to, and predictions based on those estimates would be inaccurate.  

Both of these assumptions, however, stand in sharp contrast with decades of data from psychology and neuroscience showing that repeated choices in stable environments can lead to automatic decisions driven by habits, which make people repeat previous choices rather than evaluating deliberately the relevant attributes \cite{adamsDickinson1981, dickinson1985, dolanDayan2013, perezDickinson2020, pool2022}. Under this habit mode, changes in an otherwise stable environment can produce suboptimal decision-making \cite{hardwick2019, pool2022,henriquezJara2025}. 

As people become more familiar with a stable environment, attention progressively narrows to a subset of attributes, reducing the information considered in each choice and producing choice inertia---the tendency to repeat previous choices even when doing so is no longer optimal \cite{bouton2020, thrailkill2018}. In the animal literature, a surprising change in the environment restores attention and returns control to goal-directed behavior \cite{daw2005, lee2014, bouton2024}.

In humans, there is emerging evidence that attention plays a similar role. London Underground users often stick to familiar routes until forced by a strike to explore better alternatives \cite{larcom2017}. Participants prone to habit-like strategies focus attention on previously rewarded options while ignoring superior new alternatives---a progressive narrowing of attention, or habitual attention, whereby participants subsample information and prioritize familiar cues \cite{henriquezJara2025}. These studies show that these behavioral strategies are predicted by individual differences in psychological traits, suggesting that the tendency to rely on habitual responses reflects stable differences in how cognitive resources such as attention are deployed.

What is less clear is whether the same process unfolds in DCEs when choices are repeated over long periods, just as real-life choices are. Standard DCEs use 10--20 trials---enough to estimate preferences, but not enough to induce the overtraining that shifts behavior from deliberative to automatic in laboratory tasks. In a longer, more ecologically valid design, participants may progressively narrow attention, leading to choice inertia---a pattern that would undermine the stable-preferences assumption. This possibility is also hard to test from choices alone: a person can repeat a choice while carefully checking the attributes, or while ignoring them entirely. The observed choice is the same in both cases; only the underlying process differs. A further difficulty is establishing the causal direction between attention and preferences. Attention may reflect stable preferences---people look at what they already value---or attention may actively shape preferences---what is attended to becomes what is valued. This is an open question in the decision-making literature, and standard designs, which expose participants to a small number of novel choices, cannot distinguish the two accounts.

To test this, we used a repeated transportation-choice experiment in which participants made 150 choices among three alternatives described by cost, time, comfort, and CO\textsubscript{2} emissions, while their eye movements were recorded. This is closer to the repeated choices people face in stable everyday environments than to a short DCE. During the first 100 trials the layout was stable; at trial 101 the attribute layout changed. This design allowed us to ask whether a stable environment narrows attention, whether lower attention leads to inertia, and whether estimated WTP changes when the environment is no longer stable.

The results support this account. We identified two groups of participants: one maintained exploratory attention throughout and showed little inertia; another progressively narrowed attention, developing a more stereotyped gaze pattern and a significant increase in choice inertia during the stable phase. Critically, when a manipulation that disrupted the stability of the layout was included, attention increased for all participants---and inertia fell for those who had become automatic. This shift in attention was accompanied by an immediate change in estimated preferences. Because the spatial manipulation was exogenous---it altered the layout of the display without changing the options or their attribute values---these preference shifts are attributable to the change in attentional allocation rather than to any change in the options themselves. This identifies the causal direction: attention shapes preferences, not the other way around. Repeated choices are not repeated measurements of a fixed preference: the stability of the environment shapes the decision process itself, and disrupting that stability disrupts the preferences the model recovers. Our results raise important questions about the design of DCEs and call for a deeper integration of psychological insights into economic models of choice.
\section*{Experimental Design}

The experiment was a discrete-choice task involving repeated decisions between three transportation modes: bus, metro, and ride-hailing. Each alternative was described by four attributes: cost, time, comfort, and CO\textsubscript{2} emissions. The study comprised 150 trials per participant, divided into two phases. In the training phase, participants completed 100 trials in a stable context. The colors and positions of rows (attribute labels) and columns (option labels) were fixed during this phase. In the testing phase, participants completed 50 trials in which the position of attribute rows was shuffled every five trials. This manipulation was designed to disrupt attentional habits and re-engage exploratory, goal-directed attention \citep{dayan2000,gottlieb2012}.

Attribute values in each trial were sampled through a standard DCE design procedure intended to ensure that no alternative dominated the others. In other words, no alternative was objectively better than the others across all attributes. The key manipulation was not the economic content of the options, but the stability of the visual environment. The stable phase allowed participants to learn the display. The disrupted phase broke that stable environment.

\begin{figure}[H]
    \centering
    \includegraphics[width=0.78\linewidth]{word_media/media/image1.png}
    \caption{Choice interface used in the discrete-choice task. Participants repeatedly chose between three transportation options described by cost, time, comfort, and CO\textsubscript{2} emissions.}
    \label{fig:framework}
\end{figure}

\section*{Model}

The main claim is that inertia is not fixed. In a stable environment, inertia can grow because people stop looking broadly across the display. We therefore model utility as a function of the attributes, an inertia term, and an interaction between inertia and attention:

\begin{equation}
U_{n,t,j} =
V(X_{n,t,j}, \beta_{n,t})
+ \beta_{\text{Inertia},0} I_{n,t,j}
+ \beta_{\text{Att}} \text{LatAttention}_{n,t} I_{n,t,j}
+ \epsilon_{n,t,j}.
\end{equation}

Here, $I_{n,t,j}$ indexes whether alternative $j$ is supported by the learned layout or by the previous choice, and $\text{LatAttention}_{n,t}$ is the participant's attention on trial $t$. The effective inertia parameter is therefore trial-specific:

\begin{equation}
\beta_{\text{Inertia},n,t} =
\beta_{\text{Inertia},0}
+ \beta_{\text{Inertia-Attention},n}
\text{LatAttention}_{n,t}.
\end{equation}

Attention is tied to eye-tracking through a measurement equation based on gaze entropy:

\begin{equation}
\text{Entropy}_{n,t} =
f(\text{LatAttention}_{n,t}) + \varepsilon_{n,t}.
\end{equation}

This structural approach matters because choices alone cannot distinguish a deliberate repeated choice from a habit caused by the stable layout. If attention is omitted, the model pushes the missing term into the error:

\begin{equation}
\hat{\epsilon}_{n,t,j} =
\beta_{\text{Att}} \text{LatAttention}_{n,t} I_{n,t,j}
+ \epsilon_{n,t,j}.
\end{equation}

Because reliance on the learned layout is related to lower attention, the omitted term is correlated with inertia:

\begin{equation}
\operatorname{Cov}(I_{n,t,j}, \hat{\epsilon}_{n,t,j}) \neq 0.
\end{equation}

Standard approaches that omit this latent attention interaction suffer from observational equivalence—they cannot distinguish a blind habit from a calculated choice. By resolving this endogeneity using eye-tracking as a structural measurement equation, our model can accurately isolate true preference parameters. Furthermore, attribute coefficients and willingness-to-pay are allowed to change dynamically over trials and shift after the layout disruption, capturing the true evolution of valuation.

\section*{Results}

\begin{figure*}[ht]
    \centering
    \includegraphics[width=\textwidth]{Panel.png}
    \caption{Overview of behavioral and model results. (\textbf{A}) Mean probability of choice repetition (inertia; pink, left axis) and mean gaze entropy (gray, right axis) across trials. Before the manipulation, inertia rises as entropy falls; after the layout change at trial 101, the trend reverses (teal). (\textbf{B}) Mean response time in laboratory and online samples. Both show steep RT decline during training and a partial recovery at the manipulation. (\textbf{C}) Positive correlation between log trial duration and gaze entropy: faster choices are associated with narrower attention. (\textbf{D}) Evolution of model-estimated attention, inertia ($\beta_{\text{inertia}}$), value of travel time savings (VTTS), and willingness-to-pay for CO\textsubscript{2} and comfort across trials. All parameters shift simultaneously at trial 101.}
    \label{fig:results}
\end{figure*}

We organize the results around the proposed chain: stable environment, lower attention, habit, and inertia. First, descriptive eye-tracking results show that attention narrows during training and changes after the layout is disrupted. Second, the model shows that as attention falls, inertia rises for most participants. Third, a median split on attention shows that time alone does not create the habit; inertia grows mainly when attention is low. Fourth, the estimated attribute weights and willingness-to-pay trajectories show that the stable environment changes the preferences recovered by the model.

\subsection*{Attention Narrows in the Stable Environment}

The descriptive eye-tracking results show the basic attention pattern. During the stable phase, gaze entropy decreased across trials. This means that participants looked across the display in a narrower and more regular way. After the layout changed, this pattern partly reversed, suggesting that the change made participants search again.

\begin{table}[H]
\centering
\caption{Fixed-effect model predicting entropy from trial and manipulation.}
\label{tab:entropy-model}
\begin{tabular}{lrrr}
\toprule
Predictor & Estimate & Std. Error & $t$ value \\
\midrule
Intercept & 2.070 & 0.042 & 49.786 \\
Trial & -0.004 & 0.000 & -28.264 \\
Manipulation & -0.003 & 0.046 & -0.067 \\
Trial $\times$ Manipulation & 0.002 & 0.000 & 4.347 \\
\bottomrule
\end{tabular}
\end{table}

\noindent Model: entropy $\sim$ Trial $\times$ (Trial $>$ 100) + (1 $|$ participant), $R^2 = .586$.

\subsection*{Attribute-Level Fixation Patterns Changed Across Training}

Fixation-time ratios varied by attribute and phase. During training, participants did not look at all attributes equally. After the layout changed, fixation patterns changed as well. We treat these patterns as descriptive evidence that the stable environment changed how people looked at the task.

\begin{figure}[H]
    \centering
    \includegraphics[width=0.95\linewidth]{word_media/media/image4.png}
    \caption{Fixation time ratio by attribute across 50-trial blocks. Fixation time ratio is fixation time divided by total trial time.}
    \label{fig:fixation-ratio}
\end{figure}

\begin{table}[H]
\centering
\caption{Model predicting fixation-time ratio by stage and attribute.}
\label{tab:fixation-ratio}
\begin{tabular}{lrrr}
\toprule
Predictor & Estimate & Std. Error & $t$ value \\
\midrule
Intercept & 0.303 & 0.003 & 100.182 \\
Stage 51--100 & 0.001 & 0.004 & 0.228 \\
Stage 101--150 & 0.025 & 0.005 & 5.200 \\
Time & -0.014 & 0.004 & -3.379 \\
Comfort & -0.116 & 0.004 & -26.995 \\
CO\textsubscript{2} & -0.083 & 0.004 & -19.446 \\
Stage 51--100 $\times$ Time & 0.016 & 0.006 & 2.559 \\
Stage 101--150 $\times$ Time & -0.032 & 0.007 & -4.685 \\
Stage 51--100 $\times$ Comfort & -0.013 & 0.006 & -2.104 \\
Stage 101--150 $\times$ Comfort & -0.012 & 0.007 & -1.779 \\
Stage 51--100 $\times$ CO\textsubscript{2} & -0.007 & 0.006 & -1.099 \\
Stage 101--150 $\times$ CO\textsubscript{2} & -0.056 & 0.007 & -8.243 \\
\bottomrule
\end{tabular}
\end{table}

\subsection*{Lower Attention Is Linked to More Inertia}

The model estimates an individual-level link between attention and inertia. The distribution of $\beta_{\text{Inertia-Attention}}$ is concentrated below zero (Figure~\ref{fig:histogram}). This means that, for most participants, lower attention is linked to higher inertia. A small number of participants show the opposite pattern, which suggests that people differ in how they respond to the stable environment.

\begin{figure}[H]
    \centering
    \includegraphics[width=0.95\linewidth]{histogram.jpeg}
    \caption{Distribution of the individual-level relation between attention and inertia. Negative values mean that lower attention is linked to higher inertia.}
    \label{fig:histogram}
\end{figure}

\subsection*{Time Alone Does Not Create Inertia}

Splitting trials by the median of attention shows that inertia is not just an effect of trial number. High-attention observations remain mostly flat across the task. Low-attention observations show a clear increase in inertia during the stable phase and remain elevated after the disruption (Figure~\ref{fig:median-split}). This is the key pattern: the habit appears when attention drops.

\begin{figure}[H]
    \centering
    \includegraphics[width=0.95\linewidth]{median_split_inertia.jpeg}
    \caption{Inertia over trials split by attention. High-attention choices remain mostly stable, whereas low-attention choices show a strong increase in inertia during the stable phase.}
    \label{fig:median-split}
\end{figure}

\subsection*{The Stable Environment Changes Estimated Preferences}

The beta and willingness-to-pay trajectories show that the change in attention matters for the choices themselves. During the stable phase, attention declines and inertia rises. The model also shows changes in the weights for cost, time, comfort, and CO\textsubscript{2}. Around trial 101, when the layout changes, these trajectories shift (Figure~\ref{fig:beta-evolution}). WTP for CO\textsubscript{2} and comfort changes sharply after the layout is scrambled. This suggests that a stable environment can change which attributes people use, and that breaking the stable environment can change them back.

\begin{figure}[H]
    \centering
    \includegraphics[width=0.98\linewidth]{beta_evolution.jpeg}
    \caption{Estimated trajectories for attention, attribute weights, inertia, and willingness-to-pay. The shift around trial 101 corresponds to the layout change.}
    \label{fig:beta-evolution}
\end{figure}

\section*{Discussion}

Our results show that repeated choices in a stable environment are not just repeated measurements of a stable preference. When the layout of the task is predictable, people learn to stop looking. They offload the cognitive effort of evaluating tradeoffs and rely on spatial habits. When the environment changes, attention returns, and the choices change. This has three major implications for how we understand and model decision-making.

First, choice inertia is not a static baseline; it is endogenous to attention. Standard discrete choice models often add an "inertia" dummy variable to capture why people repeat choices. But our analysis shows that time alone does not create a habit. For participants whose visual attention remained high, inertia stayed completely flat across 100 trials. Inertia only exploded for the group whose attention dropped. This means that a repeated choice is not a single phenomenon. A person might repeat a choice because they actively evaluated the attributes and still prefer it, or they might repeat it blindly because the stable layout allowed them to use a spatial shortcut. 

Second, without eye-tracking, these two behaviors look exactly the same. This is a severe case of observational equivalence. Because the decision to use a spatial heuristic is highly correlated with dropping attention, standard models that omit attention force this behavior into the error term. In our own baseline model without eye-tracking, the aggregate inertia parameter washed out and lost statistical significance. It was only by using gaze entropy to measure latent attention that we could break the tie. Eye-tracking is not just an interesting behavioral add-on; in stable environments, it is necessary to solve a fundamental endogeneity problem in the choice model.

Finally, our findings suggest that attention does more than just filter information; it dynamically constructs preferences. This provides leverage on a classic problem in neuroeconomics: do people choose an option because they look at it, or do they look at it because they already prefer it? Models like the attentional Drift Diffusion Model (aDDM) \cite{krajbich2010, krajbichRangel2011} have clearly shown that where people look predicts what they choose, but the causal arrow is often hard to pin down in static tasks. Our experimental design directly addresses this. At trial 101, the underlying economic options and their attribute levels did not change; only the spatial layout on the screen changed. There is no rational economic reason for a person's underlying preference for public transport or their willingness-to-pay for comfort to spontaneously shift at that exact moment. The only thing the manipulation forced was a return to active visual search. Because estimated willingness-to-pay shifted simultaneously with this exogenous shock to visual attention, we can isolate the causal effect. The environment forced the attentional shift, and that shift drove the change in the choice process. This supports the idea that willingness-to-pay is not a fixed number waiting to be uncovered. It is constructed in the moment, and it depends heavily on how the environment forces us to deploy our attention.

These findings connect the psychology of habits \cite{bouton2020, pool2022} with structural econometrics. For applied researchers, the lesson is practical. Long, repeated discrete choice experiments with stable layouts are not neutral measurement tools. The stable environment itself trains the participant to stop paying attention, creating the very inertia the analyst later tries to measure. If we want to capture true preferences, we have to account for the cognitive state of the person making the choice.


\matmethods{
\textbf{Participants and Experimental Task.} Thirty-nine participants completed 150 discrete choice tasks simulating a daily commute. On each trial, participants chose between three transport alternatives described by four attributes: cost, travel time, comfort, and CO\textsubscript{2} emissions. Attribute levels were fully combined and optimized using the ChoiceDesign Python package to minimize dominated alternatives and ensure an active processing baseline. To induce and test spatial habituation, the visual layout of the attributes remained perfectly stable for the first 100 trials. At trial 101, the spatial location of the attributes was unexpectedly disrupted, though the underlying economic options remained identical. The experiment was implemented in PsychoPy and conducted under controlled laboratory conditions.

\textbf{Eye-Tracking and Econometric Analysis.} Visual attention was recorded using Tobii Pro Spark eye-trackers sampling at 60 Hz. Fixation durations were computed by summing gaze events within 19 predefined Regions of Interest (ROIs), allowing us to calculate trial-specific gaze entropy and attribute-specific dwell times. After excluding trials with invalid timestamps or recording errors, the final dataset comprised 5,848 valid choice observations. Choice data and visual entropy were jointly estimated using a structural Latent Attention discrete choice model in R \cite{rstudio2020}. This model treats choice inertia as an endogenous function of visual search, using gaze entropy as a structural measurement equation to resolve observational equivalence. Full details regarding stimulus generation, eye-tracking calibration, ROI definitions, and the structural equation derivations are provided in the Supporting Information (SI) Appendix.

}

\showmatmethods{} % Display the Materials and Methods section

\dataavail{
\section{Data and Materials Availability}

All data, study materials, and analysis code used in this study are openly available at the following OSF repository: \url{https://osf.io/6xg37/}. This includes raw data, preprocessing scripts, experimental materials, and R code for the statistical analyses reported in the paper. These resources are provided to support transparency and reproducibility.

\section{Preregistration}

This study was not preregistered.

}

\acknow{
\section{Author Note}

Correspondence regarding this article should be sent to Omar D. Perez\\
(\href{mailto:omar.perez.r@uchile.cl}{omar.perez.r@uchile.cl}) or Bastian Henriquez-Jara\\
(\href{mailto:bastian.henriquez@ug.uchile.cl}{bastian.henriquez@ug.uchile.cl}).

Omar David Perez is supported by Instituto de Sistemas Complejos de Ingenieria (ISCI) ANID PIA/PUENTE AFB220002, ANID-SIA 85220023, and ANID-FONDECYT 1231027.

\section{Author Contributions}

ODP and BH designed the study. SM programmed the tasks and conducted the experiments. SM and BH analyzed the data. ODP drafted the manuscript, with critical edits by all authors.

}

\showacknow{} % Display the acknowledgments section

% Bibliography
\bibliography{bib/references}

\end{document}
